\section{Background and Motivation}
\label{sec:background}

\subsection{The \tcompile Stack}
\label{sec:stack}

\tcompile wraps a Python callable and defers compilation to the first call~\cite{DBLP:conf/asplos/AnselYHGJVBBBBC24}. Four layers then
transform the program. \emph{TorchDynamo} symbolically evaluates CPython bytecode, models Python objects with
variable trackers, extracts an FX graph~\cite{DBLP:conf/mlsys/ReedDHUA22} of tensor operations, and installs \emph{guards}: predicates
over input metadata, Python values, and global state under which the captured graph may be reused. Code that
Dynamo cannot model causes a graph break and falls back to the interpreter. \emph{AOTAutograd} traces the graph
with \emph{fake tensors} that carry only metadata, removes mutation and aliasing through functionalization,
and records a joint forward/backward graph. A \emph{decomposition} table rewrites several hundred ATen operators
into a smaller primitive set, and \emph{meta kernels} compute output shapes, dtypes, and strides without data.
\emph{TorchInductor} lowers the primitives to a loop-level IR, fuses them, and generates C++/OpenMP code for CPUs
or Triton kernels~\cite{DBLP:conf/pldi/TilletKC19} for GPUs. Compiled code is kept in an in-memory cache keyed by guards and in
on-disk caches keyed by graph and configuration hashes; the same graph can also be exported ahead of time and
packaged by AOTInductor.

Each layer can be selected in isolation through the \code{backend} argument, which \tool exploits for
localization: \code{backend="eager"} executes the graph captured by Dynamo with eager kernels (\EO),
\code{"aot\_eager"} adds AOTAutograd and decompositions (\ET), and \code{"inductor"} adds lowering and code
generation (\ER). Plain eager execution is \EZ.

\subsection{Source--Artifact Consistency}
\label{sec:definition}

Let $f$ be a Python function, $x$ an input, and $c$ an execution context that fixes tensor metadata, Python-level
values, global state, and compiler options. $E(f,x,c)$ denotes eager execution and $C(f,x,c)$ compiled execution.
We compare executions through the observation
\[
\mathit{Obs}(r)=\langle \mathit{value},\ \mathit{metadata},\ \mathit{identity},\ \mathit{exception},\ \mathit{state},\ \mathit{gradient},\ \mathit{process}\rangle ,
\]
where metadata covers shape, dtype, stride, device, and layout; identity records whether an output is an input
object or shares its storage; exception covers the type of a raised exception and whether a handler in the
program received it; state is the post-state of inputs, module buffers, and Python objects the program can
reach; and process records whether the interpreter survived. The compiler is \emph{consistent} on $(f,x,c)$ if
$\mathit{Obs}(E(f,x,c)) \approx \mathit{Obs}(C(f,x,c))$, where $\approx$ is equality except for a dtype-aware
floating-point tolerance on values and gradients. For caches we additionally compare a cold compilation in
context $c_2$ with a warm execution that first compiles under $c_1$:
$C_{\mathit{cold}}(f,x,c_2)$ versus $C_{\mathit{warm}}(f,x,c_1\!\rightarrow\!c_2)$. If eager and cold agree but warm
differs, the failure is attributed to specialization or cache reuse rather than to code generation. The context
$c$ has two parts: the program-side part (metadata, values, state) that the program can read, and the
\emph{compilation path} (graph-break positions, backend layer, execution context, artifact route) that the
program cannot observe and that must therefore never change $\mathit{Obs}$. This paper varies both.

\subsection{Motivating Examples}
\label{sec:motivation}

%% Figure: two motivating defects, each with its layer outcome strip.
\newcommand{\layercell}[3]{% #1 label, #2 colour, #3 text
  \node[draw=black!40, fill=#2, minimum width=1.5cm, minimum height=0.62cm, font=\scriptsize, align=center,
        inner sep=1pt] (#1) {#3};}
\begin{figure}[t]
\centering
\begin{minipage}[t]{0.485\linewidth}
\begin{lstlisting}
import random, torch

def augment(x):
    idx = [0, 1, 2, 3]
    random.shuffle(idx)  # Python state
    return x[idx]

g = torch.compile(augment)
x = torch.arange(4.)
g(x); g(x); g(x)   # same order each time
\end{lstlisting}
\centering
\begin{tikzpicture}[node distance=0pt]
  \layercell{a0}{tolteal!18}{\EZ eager\\varies}
  \node[right=2pt of a0] (s1) {};
  \begin{scope}[shift={(1.62,0)}]\layercell{a1}{tolred!22}{\EO Dynamo\\frozen}\end{scope}
  \begin{scope}[shift={(3.24,0)}]\layercell{a2}{tolred!22}{\ET AOT\\frozen}\end{scope}
  \begin{scope}[shift={(4.86,0)}]\layercell{a3}{tolred!22}{\ER Inductor\\frozen}\end{scope}
\end{tikzpicture}
\subcaption{\code{random.shuffle} is evaluated once at trace time: divergence enters at graph capture and is
visible only when the artifact is reused.}
\end{minipage}\hfill
\begin{minipage}[t]{0.485\linewidth}
\begin{lstlisting}
import torch

def stats(x):
    return torch.var_mean(x, correction=1)

x = torch.randn(0)        # shape[0] = 0
stats(x)                  # (nan, nan)
torch.compile(stats)(x)   # (nan, 0.)
# same for std_mean, float16, (0, 1),
# dynamic=True, and AOTInductor
\end{lstlisting}
\centering
\begin{tikzpicture}[node distance=0pt]
  \layercell{b0}{tolteal!18}{\EZ eager\\(nan, nan)}
  \begin{scope}[shift={(1.62,0)}]\layercell{b1}{tolteal!18}{\EO Dynamo\\(nan, nan)}\end{scope}
  \begin{scope}[shift={(3.24,0)}]\layercell{b2}{tolteal!18}{\ET AOT\\(nan, nan)}\end{scope}
  \begin{scope}[shift={(4.86,0)}]\layercell{b3}{tolred!22}{\ER Inductor\\(nan, 0.)}\end{scope}
\end{tikzpicture}
\subcaption{A boundary value of \code{x.shape[0]} makes Inductor fabricate a mean of zero: $E_0{=}E_1{=}E_2{\neq}E_3$
localizes the defect to code generation.}
\end{minipage}
\caption{Two source--artifact consistency failures found by \tool in PyTorch~2.14. Each strip shows the outcome
on eager execution and on the three compiled layers used for localization.}
\label{fig:motivation}
\end{figure}


Figure~\ref{fig:motivation} shows two defects found by \tool in PyTorch~2.14 that illustrate why neither graph
diversity nor an output-only oracle is sufficient.

\paragraph{A Python-semantics failure.} In Figure~\ref{fig:motivation}a, eager execution advances the global
random state, so consecutive calls return different permutations. Under \tcompile every call returns the
permutation drawn during tracing. The outcome is identical for \EO, \ET, and \ER, which places the divergence
in graph capture: Dynamo evaluates \code{random.shuffle} at trace time and embeds the result as a constant. No
tensor operator is wrong, no exception is raised, and a single-call differential test passes; the defect is
visible only when the same artifact is \emph{reused}, which is the sequence dimension of our test model.

\paragraph{A code-generation failure on a boundary shape.} In Figure~\ref{fig:motivation}b,
\code{torch.var\_mean} on an empty tensor returns \code{(nan, nan)} in eager mode and under \EO and \ET, but
\code{(nan, 0.)} under Inductor. The mean component is a fabricated zero. The trigger is a boundary value of the
factor \code{x.shape[0]}, and the layer pattern $E_0{=}E_1{=}E_2{\neq}E_3$ assigns the defect to Inductor's
reduction code. The operator is covered by PyTorch's own test suite, but not with this shape under the compiled
path.

\paragraph{Observations.} Both failures arise where a \emph{legal but uncommon context} meets a compilation
path, both are invisible to a generator that varies only graph topology, and they differ in layer, observable,
and trigger, which calls for a systematic model of context factors instead of one specialized oracle.
